% !PRINTED 668
\weberpage{668}
2. \gesperrt{By the product $\mathfrak{A}\mathfrak{B}$ of two polygons one understands the polygon that contains the points of $\mathfrak{A}$ and $\mathfrak{B}$ together, and indeed contains a point $\mathfrak{P}$ that occurs $r$ times in $\mathfrak{A}$ and $s$ times in $\mathfrak{B}$, $r+s$ times. The order of a product is equal to the sum of the orders of the factors.}

If, then, we denote distinct points by $\mathfrak{P}_1, \mathfrak{P}_2, \ldots$, we can put every polygon $\mathfrak{A}$ in one way into the form:

\weq{1}{\mathfrak{A} = \mathfrak{P}_1^{r_1}\mathfrak{P}_2^{r_2}\cdots}

The points are therefore the \gesperrt{prime elements} in the calculation with polygons. In order to have a representative for the unit as well, the \glqq zero-gon\grqq{} $\mathfrak{O}$, which contains no point at all, must be named along with them.

\gesperrt{The greatest common divisor of two polygons} $\mathfrak{A}, \mathfrak{B}$ contains every point that occurs in $\mathfrak{A}$ and $\mathfrak{B}$, and indeed as often as it occurs in each of the two.

\gesperrt{The least common multiple} contains every point of $\mathfrak{A}$ and $\mathfrak{B}$, and indeed in the highest of the orders in which it occurs in $\mathfrak{A}$ and $\mathfrak{B}$.

3. \gesperrt{If $z$ is a variable of the field $\Omega$ and $n$ the degree of the field with respect to $z$, then $z$ takes every value $c$ at $n$ points.}

For if one decomposes $z - c$ into its prime factors [§ 183, (17)]:
\[
z - c = \pi_1^{e_1}\pi_2^{e_2}\cdots, \quad e_1 + e_2 + \cdots = n,
\]
then each of these prime factors generates a point, e.\,g. $\pi_1$ the point $\mathfrak{P}_1$, at which the function $z - c$ becomes zero in the $e_1$th order. If one counts this point $e_1$-fold, one obtains a polygon $\mathfrak{P}_1^{e_1}\mathfrak{P}_2^{e_2}\cdots$ of order $n$, at whose points the function $z - c$ vanishes.

But the same theorem holds also for the value $c = \infty$, as one concludes from the fact that the degree of the field $\Omega$ with respect to $1 : z$ is the same as with respect to $z$.

4. \gesperrt{The $n$ points at which $z$ takes a value $c$ are called conjugate with respect to $z$, and the values $\eta_1, \eta_2, \ldots, \eta_n$ that a function $\eta$ in $\Omega$ takes at $n$ conjugate points are likewise called conjugate with respect to $z$, and $z$ is said to be of order $n$. The constants, and only these, have the order 0.}
